\documentclass{article}
\usepackage{graphicx}
\usepackage[table]{xcolor}
\usepackage{caption}
\usepackage{float}
\usepackage[margin=1in]{geometry}
\usepackage[colorlinks=true,linkcolor=blue]{hyperref}

\definecolor{myPurple}{RGB}{200,20,200}
\colorlet{myOrange}{red!60!yellow}
\definecolor{boxGreen}{RGB}{0,180,0}

\title{Assignment 4 - Lists, Colors, Figures etc.\\TW Lab 2025}
\author{Abha singh}
\date{\today}

\begin{document}

\maketitle

\noindent
Write \LaTeX\ code to output the following sections as it is.
Use \LaTeX\ reference and appropriate packages wherever required.

\section{Lists}

Following are some technologies that can be used to format content:

\begin{itemize}
\item \LaTeX
\item HTML (+CSS)
\item Markdown
\end{itemize}
You would most likely use \LaTeX\ in the following scenarios:
\begin{enumerate}
\item To write research papers
\item To embed math equations within other software
\item To create your CV, project reports \& presentations
\end{enumerate}
There are three main categories of electronic markup:
\begin{description}
\item[Presentational] Used by traditional word-processing systems:
binary codes embedded within document text that produce the WYSIWYG
effect. e.g. Microsoft Word
\item[Procedural] Markup is embedded in text which provides instructions
for programs to process the text. e.g. \TeX
\item[Descriptive] Markup is used to label parts of the document for
what they are. e.g. \LaTeX, HTML5
\end{description}
\vspace{7cm}
\section{Tables \& Colors}
Code for \ref{mytable} is written in this subsection. However, it will
be rendered on a special page. Change the optional arguments such that the table is drawn at the
\fcolorbox{boxGreen}{white}{top} or
\fcolorbox{boxGreen}{white}{bottom} of a page or
\fcolorbox{boxGreen}{white}{within} this section. The last row of this table shows the width
of each column.
\smallskip
\colorbox{orange}{Read more about how colors are represented as numbers.}
Try to find out about the HTML notation. Define a custom color with 200 red, 20 green and 200 blue. Define another custom color using
\texttt{\textbackslash colorlet} with 60\% red and 40\% yellow.

\[
\textcolor{myPurple}{T(n)}
\textcolor{red}{=}
\textcolor{myOrange}{aT(n/b)}
\textcolor{red}{+}
\textcolor{myOrange}{f(n)}
\]

\section{Figures}

\subsection{Pictures: rotation \& cropping}

Include the MNIT logo and rotate it appropriately. Also include the MNIT building picture but crop the sky and make it as wide as the text.

\begin{figure}[h]
\centering
\includegraphics[width=0.3\textwidth,angle=-60]{logo.png}
\caption{MNIT logo}
\end{figure}

\begin{figure}[h]
\centering
\includegraphics[
trim=0pt 0pt 0pt 40pt,
clip,
width=1\textwidth
]{building.png}
\caption{This is a floating image and can appear anywhere.}
\end{figure}

\vspace{4cm}

\subsection{Multiple images}

This subsection includes two minipages with equal width.

\begin{figure}[h]
\centering

\begin{minipage}{0.495\textwidth}
\centering
\includegraphics[
width=0.8\linewidth,
]{logo.png}
\caption{This is the first minipage}
\end{minipage}
\hfill
\begin{minipage}{0.495\textwidth}
\centering
\includegraphics[
width=0.8\linewidth
]{building.png}
\caption{This is the second minipage}
\end{minipage}

\end{figure}

\newpage
\begin{table}[t]
\centering

\begin{tabular}{|p{2.5cm}|p{2cm}|p{2cm}|p{3.5cm}|}
\hline
\rowcolor{green}
Language & Requires compilation & Output & Use \\
\hline
C/C++ & Yes & Executable & To make apps \\
\hline
\LaTeX & Yes & PDF & To make documents \\
\hline
HTML/CSS/JS & No & Webpage & To make websites \\
\hline
\rowcolor{yellow}
2.5 cm & 2 cm & 2 cm & 3.5 cm \\
\hline
\end{tabular}

\caption{This table should be rendered on a special page.}
\label{mytable}

\end{table}

\end{document}


